\documentclass[11pt]{article}
\usepackage[margin=1in]{geometry}
\usepackage{amsmath,amssymb,amsthm,hyperref}
\setlength{\emergencystretch}{2em}
\newtheorem{theorem}{Theorem}
\title{An observable can miss the slow mixing mode\\Disproof of Conjecture 00000003429}
\author{ziangni-sys}
\date{October 8, 2026}
\begin{document}
\maketitle

\paragraph{Scope.}
Both source language versions identify the autocorrelation decay of
observables with the mixing spectral radius and propose lag regression
to estimate that radius. An individual observable need not detect the
slowest nonconstant mode. We give a strictly positive, lazy, reversible
finite-state chain and a nonzero centered observable whose exact
exponential decay factor is strictly smaller than the mixing spectral
radius. The example refutes the blanket equality for observables.
It does not refute a worst-case upper bound or an equality obtained by
taking the supremum over all observables, nor a claim restricted to
observables having nonzero projection onto the slowest mode.

\paragraph{A genuine Markov chain.}
On the state space $\{0,1,2\}$, use the transition matrix
\[
 P=\begin{pmatrix}
 5/8&1/8&1/4\\
 1/8&5/8&1/4\\
 1/4&1/4&1/2
 \end{pmatrix},\qquad \pi=(1/3,1/3,1/3).
\]
Every entry is strictly positive and each row sums to one. Thus the chain
is irreducible and aperiodic. Its diagonal entries are at least $1/2$, so
it is lazy. Symmetry gives both stationarity $\pi P=\pi$ and detailed
balance $\pi_iP_{ij}=\pi_jP_{ji}$.

\paragraph{The actual centered restriction.}
The centered space is
$H_0=\{v:\sum_i\pi_i v_i=0\}$.
Set $g=(1,-1,0)$ and $f=(1,1,-2)$. These vectors form a basis of $H_0$,
and direct multiplication gives
\[
                   Pg=\tfrac12g,\qquad Pf=\tfrac14f.
\]
For completeness, the coordinate map and its inverse are explicit:
\[
 E(a,b)=(a+b,-a+b,-2b),\qquad
 a=\frac{v_0-v_1}{2},\quad b=\frac{v_0+v_1}{2}
       \quad(v\in H_0).
\]
Thus $E$ is a linear bijection onto $H_0$ and
\[
        PE=ED,\qquad D=\begin{pmatrix}1/2&0\\0&1/4\end{pmatrix}.
\]
The spectrum of this restriction is exactly $\{1/2,1/4\}$, as also follows
from $\det(zI-D)=(z-1/2)(z-1/4)$. The same diagonalization holds on
complexifying $H_0$, so no nonreal spectral values arise. The mixing
spectral radius, the largest modulus on the centered restriction, is
therefore
\[
                         \rho_{\mathrm{mix}}=\tfrac12.
\]

\begin{theorem}
For the nonzero centered observable $f=(1,1,-2)$, the normalized stationary
autocorrelation is $C_f(n)=4^{-n}$. Its exponential decay factor and every
consecutive-lag regression ratio equal $1/4$, whereas
$\rho_{\mathrm{mix}}=1/2$.
\end{theorem}
\begin{proof}
Induction gives $P^nf=4^{-n}f$. In stationarity, the Markov property gives
\[
 \operatorname{Cov}_\pi(f(X_0),f(X_n))
   =\sum_i\pi_i f_i(P^nf)_i
   =4^{-n}\frac{1+1+4}{3}=2\cdot4^{-n},
\]
because $\sum_i\pi_i f_i=0$. In particular $\operatorname{Var}_\pi(f)=2>0$,
so normalization gives $C_f(n)=4^{-n}$. For all $n\geq0$,
\[
 \frac{C_f(n+1)}{C_f(n)}=\frac14,\qquad
 |C_f(n+1)|^{1/(n+1)}=\frac14.
\]
Hence the positive-root exponential-rate limit is exactly $1/4$,
strictly smaller than $1/2$. In logarithmic conventions the rates are
$\log4$ and $\log2$, likewise unequal.
\end{proof}

\paragraph{Why the spectral information is lost.}
Under the stationary inner product, $\langle f,g\rangle_\pi=0$.
This observable has no component in the eigenvalue-$1/2$ direction.
There is no defect in positivity, laziness, reversibility, centering, or
variance: it simply probes a different nonconstant eigenspace.
Lag regression cannot recover an unobserved mode from this autocorrelation.

\paragraph{Lean coverage.}
The project defines the actual transition matrix and stationary weights,
proves positivity, stochasticity, laziness, stationarity and detailed
balance, and uses explicit coordinates to identify the entire centered
restriction. It computes its spectrum and spectral radius via matrix
invertibility. Autocovariance is the actual stationary finite-sum formula
using powers of $P$. Its exact value, normalized correlation, all lag
ratios and the root-rate limit are proved, ending in the inequality
between the two rates. No eigenvalue, covariance identity, or decay-rate
certificate is assumed as a hypothesis.
\end{document}
